\documentclass[10pt,a4paper]{amsart}
\usepackage[margin=19mm]{geometry}
\usepackage{amsmath,amssymb,mathtools}
\usepackage[scr=rsfs,cal=euler]{mathalfa}
\usepackage{xfrac}
\usepackage[unicode,hidelinks,bookmarks=false]{hyperref}
\newcommand{\ie}{i.e.\ }
\newcommand{\cf}{cf.\ }
\newcommand{\resp}{respectively\ }
\newcommand{\Subsection}{\subsection}
\providecommand{\nobreakdash}{\nobreak\hbox{-}}
\setlength{\emergencystretch}{2em}
\multlinegap0pt
\usepackage{bm}
\usepackage{bbm}
\usepackage{dsfont}
\usepackage{xspace}
\usepackage{cancel}
\usepackage{mathtools}
\usepackage{amsthm}
\makeatletter
\@ifundefined{theorem}{\newtheorem{theorem}{Theorem}}{}
\@ifundefined{claim}{\newtheorem{claim}[theorem]{Claim}}{}
\@ifundefined{lemma}{\newtheorem{lemma}[theorem]{Lemma}}{}
\@ifundefined{proposition}{\newtheorem{proposition}[theorem]{Proposition}}{}
\makeatother
\newcommand{\R}{\mathbb R}
\newcommand{\lsm}{\lesssim}
\newcommand{\vertiii}[1]{{\left\vert\kern-0.25ex\left\vert\kern-0.25ex\left\vert #1 
    \right\vert\kern-0.25ex\right\vert\kern-0.25ex\right\vert}}
\title[The uniqueness of the ground state and the dynamics of nonli]{The uniqueness of the ground state and the dynamics of nonlinear Schr\"odinger equation with inverse square potential}
\author{Kai Yang, Chongchun Zeng, Xiaoyi Zhang}
\date{Source: arXiv:2603.10338v1 (CC BY 4.0)}
\begin{document}
\maketitle

\noindent\emph{Source and licence.} The uniqueness of the ground state and the dynamics of nonlinear Schr\"odinger equation with inverse square potential. Version arXiv:2603.10338v1 (CC BY 4.0), distributed under the Creative Commons Attribution 4.0 International licence, as recorded in the arXiv licence metadata for this exact version and shown on the abstract page https://arxiv.org/abs/2603.10338v1. Full rights notice: licenses/arxiv-2603.10338-CC-BY-4.0.txt.\par\smallskip

\noindent\textit{Editorial scope.} Counted unit: the Theorem whose source label is \texttt{exiuni} in the article (source0.tex lines 2295--2310, source label \texttt{exiuni}), delivered together with the article's own complete original proof of it (source0.tex lines 2313--2508, one \texttt{proof} environment of 196 lines). No mathematics is written, repaired or re-derived: statement and proof are the article's own text line for line. Required context retained uncounted: 1102 (lines 1996--1998); prop:linzeng (lines 2008--2046); R2 (lines 2073--2082); Lemma3.1 (lines 2196--2206); 918 (lines 2200--2205); JJ6 (lines 2263--2267). Every definition, display and label of those lines is retained unchanged. Omitted from this package and not redistributed: 1--1995, 1999--2007, 2047--2072, 2083--2195, 2206--2262, 2268--2294, 2311--2312, 2509--3997. Nothing inside a retained excerpt is omitted or altered: every retained body is delivered exactly as the article's lines are written, so no omission receipt of any kind accompanies this package. Citations: the retained excerpts carry no citation command at all, so no bibliography entry of the article is transcribed and no reference is redistributed. Reference adaptation: none was needed and none was made; every reference inside the delivered unit resolves inside it, so no unresolved reference or citation is produced. Printed slips kept literal: no printed word, symbol, hypothesis or number of the counted statement or of its proof was altered. Layout and class: the article's own class is \texttt{amsart} with packages including amssymb, amsmath, eurosym, paralist, graphics, epsfig, tabu, graphicx, epstopdf, hyperref, amsfonts, fancyhdr, geometry; the wrapper is the lane's pinned \texttt{amsart} editorial wrapper at 10pt on A4, which restates the theorem-like environments the retained text needs and adds the article's own macro definitions that the retained text needs (\texttt{\textbackslash R}, \texttt{\textbackslash lsm}, \texttt{\textbackslash vertiii}), with extra wrapper packages (bm, bbm, dsfont, xspace, cancel, mathtools). The delivered numbering is the unit's own. Assets: no retained span contains a figure environment, a graphics-inclusion command, a PDF-page inclusion, a table or a TikZ picture; no image file is copied into this package, the package declares no asset dependency and no image file is redistributed with it. Licence: version arXiv:2603.10338v1 of this article is distributed under the Creative Commons Attribution 4.0 International licence, as recorded in the arXiv licence metadata for this exact version and shown on the abstract page https://arxiv.org/abs/2603.10338v1. A full rights notice is delivered at licenses/arxiv-2603.10338-CC-BY-4.0.txt. Characters: every retained excerpt is pure ASCII, so the delivered engine, XeLaTeX with its default Latin Modern text font, reports no missing character.
\begin{align}\label{1102}
v_t=JL v+R(v),
\end{align}

\begin{proposition}\label{prop:linzeng}
The flow $e^{tJL}$ is a well-defined group of bounded linear operators on $(H^1(\R^d))^2$ and there exist closed subspaces $E^u,\ E^s$ and $E^c$ such that 

a) 
%$H^1 = E^u \oplus E^s \oplus E^c$ and 
$\dim E^u=\dim E^s=1$.

b) $e^{tJL}(E^{u,s,c})=E^{u,s,c}$.

c) $ \langle Lu, u\rangle =0, \; \forall u\in E^{u,s},$ and $$E^c=\{u\in (H^1(\R^d))^2 \mid \langle Lu, v\rangle=0, \ \forall v\in E^u\oplus E^s\}.$$

%d) There exist $c,\lambda>0$ such that 
%$$
%JLu=\pm \lambda u, \forall u\in E^{u,s}, \mbox{ hence }\biggl |e^{tJL}|_{E^{u,s}}\biggr|\le Ce^{\mp \lambda t}, \ \forall t\ge 0. 
%$$

d) $\bigl |e^{tJL}|_{E^c}\big|\le C(1+|t|)$, $\forall t\in \mathbb R $.

e) There exist $\epsilon >0$, $e_0>0$, and a closed subspace $E^e$ such that $E^c=\ker L\oplus E^e$ and $(H^1(\R^d))^2=E^u\oplus E^s\oplus \ker L \oplus E^e$. Moreover,  

$$L\sim 
\begin{pmatrix} 0 & 1 & 0 & 0\\
1&0&0&0\\
0&0&0&0\\
0&0&0&L_e
\end{pmatrix}, \quad 
JL\sim \begin{pmatrix}
e_0 &0&0&0\\
0&-e_0 &0&0\\
0&0&0&A_{0e}\\
0&0&0&A_e
\end{pmatrix}
$$
and 
$$
L_e\ge \epsilon>0, \quad \langle L_e e^{tA_e}u, e^{tA_e}v\rangle=\langle L_e u, v\rangle.
$$   
Here $\sim$ indicates that the operators $L$ and $JL$ take those block forms in the direct sum decomposition. $L_e: E^e\to (E^e)^*$ is defined as $\langle L_e u, v\rangle=\langle Lu, v\rangle$, for any $u, v\in E^e$. Operators $A_e$ and $A_{0e}$ are defined as the projections of $JL|_{E^c}$ onto subspaces $E^e$ and   $\ker L$ based on the above direct sum decomposition. In particular, $A_{0e} \in L(E^e, \ker L)$.      
\end{proposition}

\begin{align}\label{R2}
\begin{cases}
|R(v)|\lsm Q^{p-1}|v|^2+ |v|^{p+1}, \qquad |R(v_1)-R(v_2)|\lsm |v_1-v_2|\sum_{j=1}^2(Q^{p-1}|v_j|+|v_j|^p), \\
|\nabla R(v)|\lsm (Q^{p-1}|v|+|v|^p) \bigl(|\nabla v|+|Q^{-1}\nabla Q v|\bigr),\\
|\nabla R(v_1)-\nabla R(v_2)|\lsm  \bigl(|Q^{-1}\nabla Q (v_1-v_2)|+|\nabla (v_1-v_2)\bigr)\sum_{j=1}^2(Q^{p-1}|v_j|+ |v_j|^p)\\
\qquad \qquad\qquad+|v_1-v_2|\big(\sum_{j=1}^2 (|\nabla v_j|  +|Q^{-1}\nabla Q v_j|)\big)
%+|Q^{-1}\nabla Q v_i|\bigr)
\bigl(Q^{p-1}+\sum_{j=1}^2|v_j|^{p-1}\bigr). 
\end{cases}
\end{align}

\begin{lemma}\label{Lemma3.1}
	Let $u_0\in E^c$, $f(t)\in E^c$ and $T>0$. Let 
\[u(t,x)=e^{tJL} u_0, \quad v(t,x)=\int_0^t e^{(t-s)JL}f(s)ds,\quad w(t,x)=\int_t^T e^{(t-s)JL}f(s)ds.\]
Then,  
	\begin{align}
		\|u(\pm t)\|_{ H^1} &\lesssim \left\langle t \right\rangle \|u_0\|_{ H^1}, \label{H1bd}\\
		% \|u(\pm t)\|_2&\lsm \|u_0\|_2+\left\langle t\right\rangle^2\|u_0\|_{\dot H^1},\label{607}\\
		\|u\|_{ S^1([0,T])} &\lesssim \left\langle T \right \rangle^2 \|u_0\|_{ H^1}. \label{local u}\\
		\|v\|_{ S^1([0,T])}+\|w\|_{ S^1([0,T])}&\lsm \langle T\rangle ^2 \|f\|_{L_t^1 H^1([0,T])}. \label{918}
	\end{align}
\end{lemma}

	\begin{align}
		\|v\|_{ S^1(I_j)} &\lesssim \|v(j\eta)\|_{H^1} + \|v-\bigl((p+1)Q^p v_1 +iQ^pv_2\bigr)\|_{L_t^1 H^1(I_j)}
                  +\|f\|_{L_t^1H^1(I_j)}\label{JJ6} \\
		&\lesssim \|v(j\eta)\|_{H^1} +\eta^{1-\frac 1{q^*}}\|v\|_{ S^1(I_j)} +\|f\|_{L_t^1H^1(I_j)}. \notag
	\end{align}

\begin{theorem}\label{exiuni} 
Let $d=3,4,5$ and $p\in (\tfrac 4d, \tfrac 4{d-2})\cap [1,\infty)$. Let $0>a>\tfrac{(d-2)^2}4(-1+\frac{p^2}{(p+1)^2}).$ There exists a constant $C>0$ such that, for any $\lambda\in(0, e_0]$ and $y_0^- \in (-\delta, \delta)$ where 
% $\delta\in(0,\delta_\lambda)$ with 
$\delta=\frac 1{C} \min(\lambda,\lambda^4)$, equation \eqref{1102} admits a unique solution satisfying  
\begin{equation} \label{E:class}
v(0)=y_0^-V^-+y_0^+V^++v^c(0), \textit{ and } \|v(t) \|_{H^1}
%\dot S^1([t,\infty)\times\R^3)}
\le C \delta e^{-\lambda t},  \quad \forall t\ge 0.
\end{equation}
Moreover, the following estimates hold:
\begin{align}\label{438}
& |y_0^+|+\|v^c(0)\|_{H^1_x}\le C |y_0^-|^2, \\
& \|v\|_{ S^1([t,\infty))}\le C^2 \delta e^{-e_0 t}, \quad \forall t\ge 0. 
\end{align}
For any $y_0$, $\tilde y_0$ with $y_0^-\tilde y_0^->0$ and $|y_0^-|, |\tilde y_0^-|< \delta$, the corresponding solutions satisfy $v(t,x)=\tilde v(t+T,x)$ for some $T=T(y_0, \tilde y_0)$. 
\end{theorem}

\begin{proof} 
Using the decomposition $(H^1(\R^d))^2=E^s\oplus E^u\oplus E^c$ from Proposition \ref{prop:linzeng}, we express: 
\begin{align}\label{decu}
v=y^+V^++y^-V^-+v^c
\end{align}
where $y^{\pm}=\langle LV^{\mp}, v\rangle$ and $v^c=v-y^-V^--y^+V^+$. Equation \eqref{1102}  decouples into:
\begin{align*}
\begin{cases}
\dot y^-&=-e_0 y^-+R^-(v)\\
\dot y^+&=e_0 y^++R^+(v)\\
\frac{\partial}{\partial t}v^c&=JL v^c+R^c(v), 
\end{cases}
\end{align*}
whose Duhamel formulation (incorporating exponential decay) becomes:
\begin{align}\label{inteqn}
\begin{cases}
y^-(t)&=e^{-e_0 t} y_0^-+\int_0^t e^{-e_0(t-s) }R^-(v(s))ds\\
y^+(t)&=-\int_t^\infty e^{e_0(t-s)}R^+(v(s))ds\\
v^c(t)&=-\int_t^\infty e^{JL(t-s)}R^c(v(s))ds. 
\end{cases}
\end{align}
Here, $R^{\pm}(v)$ and $R^c(v)$ are defined via
\begin{align*}
R^{\pm}(v)=\langle LV^{\mp}, R(v)\rangle, \ R^c(v)=R(v)-R^+(v)V^+-R^-(v)V^-. 
\end{align*}
For functions $v(t, x)$ expressed as in \eqref{decu}, define the norm 
\[
\vertiii{ v}=\max\bigl\{\sup_{t\ge 0} e^{\lambda t}|y^+(t)|, \ \sup_{t\ge 0} e^{\lambda t}|y^-(t)|,\ \sup_{t\ge 0} e^{\lambda t}\|v^c\|_{S^1([t,\infty))}\bigr\},
\]
and define the mapping: 
\[
\Phi (v) (t,x) = \tilde v (t, x)  = \tilde y^+  (t) V^+(x) + \tilde y^-(t) V^-(x) + \tilde v^c(t, x),
\]
where $\tilde y^\pm (t)$ and $\tilde v^c$ are given by the right sides of \eqref{inteqn}. 

% we and equip the space with the norm:
We shall prove $\Phi$ is a contraction on the closed ball:
%defined by 
\begin{align*}
B_{\delta,\lambda}=\bigl \{v = y^+ V^+ + y^- V^-+ v^c \bigl | y^{\pm} \in C^0([0, \infty))\;, v^c\in  S^1 ([0, \infty)),\; \vertiii{v}\le 2\delta.\bigr\}
\end{align*}
It is easy to see $B_{\delta, \lambda}$ increases in $\delta$ and decreases in $\lambda$. 

For $v_1, v_2 \in B_{\delta, \lambda}$ expressed as:  
\[v_j=y_j^+V^++y_j^-V^-+v_j^c\in B_{\delta, \lambda}, \; j=1,2. \]
we denote their images under $\Phi$ by $\tilde v_1$ and $\tilde v_2$. The contraction property follows from estimates of the nonlinear terms $R^{\pm}$ and $R^c$, which we now detail. 

\begin{claim}\label{JJ3} For $p \in [1, \frac 4{d-2})$, there exists $C>0$ such that, for any $v_1, v_2 \in B_{\delta, \lambda}$, the following estimates hold:
\begin{align*}
\begin{cases}
|R^{\pm}(v_1(s))|\le C\delta^2 e^{-2\lambda s},\\
|R^{\pm}(v_1(s))-R^{\pm}(v_2(s))|\le C \delta e^{-2\lambda s}\vertiii{v_1-v_2},\\
\|R^c(v_1)\|_{L_t^1 H^1([T_0, T_0+T_1])}\le C\langle T_1\rangle^{1-\frac 2{q^*}}\delta^2 e^{-2\lambda T_0}, \\
\|R^c(v_1)-R^c(v_2)\|_{L_t^1 H^1([T_0, T_0+T_1])}\le C\langle T_1\rangle ^{1-\frac 2{q^*}}\delta e^{-2\lambda T_0}
\vertiii{v_1-v_2}. 
\end{cases}
\end{align*}
\end{claim}
We focus on proving the difference estimates (the second and the fourth inequalities) as the rest follow by setting one of $v_i=0$. 

From the pointwise estimates \eqref{R2} and H\"older inequality, we compute: 
\begin{align*}
|R^{\pm}(v_1(s))-R^{\pm}(v_2(s))|&=|\langle LV^{\mp}, R(v_1)-R(v_2)\rangle |\\
&\lesssim  \|LV^{\mp}\|_{L_x^{\frac{4d}{2d-p(d-2)}}}\|v_1-v_2\|_{L_x^{\frac{4d}{2d-p(d-2)}}} \sum_{j=1,2}\bigl(\|Q\|_{L_x^{2^*}}^{p-1}\|v_j\|_{L_x^{2^*}}+\|v_j\|_{L_x^{2^*}}^p\bigr)
\\
&\lesssim \|v_1-v_2\|_{H^1}\sum_{j=1,2}( \|v_j\|_{H^1}+\|v_j\|_{H^1}^p)\\
&\le C \delta e^{-2\lambda s}\vertiii{v_1-v_2}. 
\end{align*}
Here we have used $\frac{4d}{2d-p(d-2)} \in (2, 2^*)$.


Similarly, for the gradient term, we have
\begin{align}
\|\nabla(R(v_1)-R(v_2))&\|_{L_t^1 L_x^2([T_0,T_0+T_1])}\le \|\nabla(v_1-v_2)\|_{L_t^{q^*}H_x^{1, r^*}}\label{525}\\
&\qquad\cdot\bigl[
T_1^{1-\frac 2{q^*}}\|Q\|_{L_x^{r}}^{p-1}\sum_{j=1,2}\|v_j\|_{L_t^{q^*}L_x^r}
+T_1^{1-\frac {p+1}{q^*}}\sum_{j=1,2}\|v_j\|_{L_t^{q^*}L_x^r}^p\bigr]\notag\\
&+\|v_1-v_2\|_{L_t^{q^*}L_x^{r}}\sum_{j=1,2} \big( \|\nabla v_j\|_{L_t^{q^*}L_x^{r^*}} +  \|v_j\|_{L_t^{q^*}L_x^{r^*}}  \big) \notag\\
&\qquad\cdot\bigl[T_1^{1-\frac 2{q^*}}\|Q\|_{L_x^{r}}^{p-1}+T_1^{1-\frac {p+1}{q^*}}
\sum_{j=1,2}\|v_j\|_{L_t^{q^*}L_x^r}^{p-1}\bigr]\notag\\
&\le C\langle T_1\rangle^{1-\frac {2}{q^*}}\delta e^{-2\lambda T_0}\vertiii{v_1-v_2}. \notag
\end{align}
The same bound holds for: 
\[
\|R(v_1)-R(v_2)\|_{L_t^1L_x^2([T_0, T_0+T_1])}\le C\langle T_1\rangle^{1-\frac {2}{q^*}}\delta e^{-2\lambda T_0}\vertiii{v_1-v_2}.
\]
These two estimates together with the boundedness of the projection onto center space yield the last estimate in Claim \ref{JJ3}.

We now verify the contraction property for $\Phi$. For any $v_1, v_2\in B_{\delta, \lambda}$, Claim \ref{JJ3} implies 
\begin{align*}
	e^{\lambda t}\bigl|\tilde y_1^-(t)-\tilde y_2^-(t)\bigr| &\le  e^{\lambda t}
	  \int_0^t e^{-e_0(t-s)}\bigl |R^-(v_1(s))-R^-(v_2(s))\bigr|ds \\
	&\le C\delta \int_0^t e^{(\lambda- e_0)(t-s)} e^{-\lambda s}ds \vertiii{v_1-v_2}\\
	&\le \frac 14\vertiii{v_1-v_2}, \; \forall t\ge0,
\end{align*}
for $\delta$ sufficiently small as in Theorem \ref{exiuni}. Similarly, 
\begin{align}
		e^{\lambda t}\bigl |\tilde y_1^+(t)-\tilde y_2^+(t)\bigr |&\le C\delta\int_t^{\infty}e^{(\lambda+e_0)(t-s)}e^{-\lambda s}ds\vertiii{v_1-v_2}\le \frac 14\vertiii{v_1-v_2}. \label{1009}
\end{align}
%for the same choice of $\delta$. 

For Strichartz estimate of $\tilde v_1^c-\tilde v_2^c$, we decompose
\begin{align*}
	\|\tilde v_1^c-&\tilde v_2^c\|_{ S^1([T,\infty))}\\
%	\biggl\|\int_t^{\infty} e^{(t-s)JL}R^c(v(s))ds\biggr \|_{\dot S^1([T,\infty))} \\
%	%&\le \sum_{N\ge 1, N\in\mathbb N
	%			}\biggl\|\int_t^{\infty} e^{(t-s)JL} R^c(v(s))ds \biggr\|_{\dot S^1([NT, (N+1)T])} \\
	&\le \sum_{
		N\ge 1, N\in\mathbb N
		} \biggl\|\int_t^{T +N}e^{(t-s)JL} \bigl(R^c(v_1(s))-R^c(v_2(s))\bigr)ds\biggr\|_{ S^1([T+ N-1, T+ N])}  \\
	&\quad +\sum_{
		N\ge 1, N\in\mathbb N
		} \biggl\|\int_{T+N}^{\infty} e^{(t-s)JL}\bigl(R^c(v_1(s))-R^c(v_2(s))\bigr)ds\biggr\|_{S^1([T+N-1, T+ N])} \\
	&:= I + II.
\end{align*}
Using \eqref{918} and Claim \ref{JJ3} yields: 
\begin{align*}
	I &\le \sum_{
		N\ge 1, N\in\mathbb N} 
		%\langle T\rangle^2 
		\|R^c(v_1(s))-R^c(v_2(s))\|_{L_t^1 H^1([T+N-1, T+N])} \\
	&\le C\delta\vertiii{v_1-v_2} \sum_{
		N\ge 1, N\in\mathbb N		} 
		%\langle T\rangle^3 
		 e^{-2\lambda (T+N-1)}\\
	& \le \frac 1{\lambda} C \delta  e^{-2\lambda T}\vertiii{v_1-v_2}\le \frac 18 \vertiii{v_1-v_2}e^{-2\lambda T}.
	 % \le \delta e^{-\lambda T}.
\end{align*}
For the term $II$, we further partition the integral and compute 
\begin{align*}
	II &\le \sum_{
		N\ge 1	}\sum_{
			 M\ge N+1
		 }\biggl \|\int_{T+M-1}^{T+M} e^{(t-s)JL} 
	     \bigl (R^c(v_1(s))-R^c(v_2(s))\bigr)ds\biggr\|_{ S^1([T+N-1, T+N])} \\
	   &\le \sum_{
			   N\ge 1,
			   M\ge N+1
               }\int_{T+M-1}^{T+M}\bigl \|e^{(t-s)JL}  \bigl (R^c(v_1(s))-R^c(v_2(s))\bigr)\bigr\|_{ S^1([T+N-1, T+N])}ds. 
\end{align*}
Note $0 \le s-t \le M+1$, applying Lemma \ref{Lemma3.1} and Claim \ref{JJ3} we obtain
\begin{align*}               
	   II&\le \sum_{
	N\ge 1, M\ge N+1        
	      }\int_{T+M-1}^{T+M} (M+1)^2\bigl \|R^c(v_1(s))-R^c(v_2(s))\bigr\|_{ H^1} ds\\
	   &\le \sum_{
			   M\ge2
                         }(M+1)^3 \|R^c(v_1(s))-R^c(v_2(s))\|_{L_t^1 H^1([T+M-1, T+M])} \\
	   &\le C\delta \vertiii{v_1-v_2} \sum_{
			   M\ge 2
                   	      }(M+1)^3  e^{-2\lambda (T+M-1)} \\
%	   &\lesssim (2\delta)^2 T^2 e^{-\frac 32 e_0 T} \\
	  & \le \frac C{\lambda^4} \delta e^{-2\lambda T}\vertiii{v_1-v_2}\le \frac 18\vertiii{v_1-v_2} e^{-2\lambda T}. 
	   \end{align*}
Combining these yields:
%and combining the case $T<1$, 
\begin{align}\label{1016}
\sup_{T\ge 0}e^{\lambda T}\|\tilde v^1_c-\tilde v^2_c\|_{ S^1([T,\infty))}\le\frac 14 \vertiii{v_1-v_2}.
\end{align}
This proves that 
\begin{align*}
\vertiii{\tilde v_1-\tilde v_2}\le \frac 12\vertiii{v_1-v_2}.
\end{align*}
Since $\vertiii {\Phi(0)} = |y_0^-| < \delta$,  the contraction estimate $Lip( \Phi) \le \frac 12$ also implies that  
%A similar argument shows 
$\Phi$ maps $B_{\delta,\lambda}$ into itself. The existence and uniqueness of solutions to \eqref{inteqn} in $B_{\delta,\lambda}$ are proved. 


  



To prove the uniqueness in the class defined by \eqref{E:class}, we verify that 
exponential decay in $H^1(\R^d)$ implies the same decay as in the Strichartz space $S^1$. Specifically, given 
%solution  
%assume $v(t)$ is a solution to \eqref{1102} satisfying   
\[
\|v(t) \|_{H^1} \lsm e^{-\lambda t}, \quad t\ge 0,
\]
we shall prove
\begin{align}\label{a101}
\| v\|_{S^1 ([t, \infty))} \lsm e^{-\lambda t}, \quad t\ge 0.
\end{align}
Indeed, for any $\tau\ge 0$ and sufficiently small $\eta>0$ (independent of $\tau$), applying Strichartz estimate on $[\tau, \tau+\eta]$ and follow the similar procedure as in \eqref{JJ6}, \eqref{525}, we obtain
\begin{align*}
\|v\|_{ S^1{([\tau,\tau+\eta])}}\le &C\|v(\tau)\|_{H^1}+C\eta^{1-\frac 1{q^*}}\|v\|_{ S^1{([\tau,\tau+\eta])}}\\
&+C\eta^{1-\frac 2{q^*}}\|v\|_{ S^1{([\tau,\tau+\eta])}}^2+C\eta^{1-\frac{p+1}{q^*}}\|v\|_{ S^1{([\tau,\tau+\eta])}}^{p+1}.
\end{align*}
 Recall $\|v(\tau)\|_{H^1}\lsm e^{-\lambda \tau}$. Choosing sufficiently small $\eta$ and using a standard argument gives 
\begin{align*}
\|v\|_{S^1([\tau,\tau+\eta])}\le 2C\|v(\tau)\|_{H^1}\lsm e^{-\lambda \tau}. 
\end{align*}
The estimate of $\|v\|_{ S^1([t,\infty))}$ then follows by partitioning the interval $[t,\infty)$ into subintervals of length $\eta$ and summing all the local estimates.

Finally, the quadratic estimate \eqref{438} follows from the existence and uniqueness of solutions in the smaller ball $B_{C |y_0^-|, e_0}$ and inequality \eqref{1009} with $y_2^+ =0$. The time translation invariance follows from the uniqueness and decay of solutions. The proof of Theorem \ref{exiuni} is complete. 
\end{proof}

\end{document}
